\documentclass[a4paper,12pt,twoside]{memoir}

% Castellano
\usepackage[spanish,es-tabla]{babel}
\selectlanguage{spanish}
\usepackage[utf8]{inputenc}
\usepackage[T1]{fontenc}
\usepackage{lmodern} % Scalable font
\usepackage{microtype}
\usepackage{placeins}

\RequirePackage{booktabs}
\RequirePackage[table]{xcolor}
\RequirePackage{xtab}
\RequirePackage{multirow}

% Links
\PassOptionsToPackage{hyphens}{url}\usepackage[colorlinks]{hyperref}
\hypersetup{
	allcolors = {red}
}

% Ecuaciones
\usepackage{amsmath}

% Rutas de fichero / paquete
\newcommand{\ruta}[1]{{\sffamily #1}}

% Párrafos
\nonzeroparskip

% Huérfanas y viudas
\widowpenalty100000
\clubpenalty100000

% Imágenes

% Comando para insertar una imagen en un lugar concreto.
% Los parámetros son:
% 1 --> Ruta absoluta/relativa de la figura
% 2 --> Texto a pie de figura
% 3 --> Tamaño en tanto por uno relativo al ancho de página
\usepackage{graphicx}
\newcommand{\imagen}[3]{
	\begin{figure}[!h]
		\centering
		\includegraphics[width=#3\textwidth]{#1}
		\caption{#2}\label{fig:#1}
	\end{figure}
	\FloatBarrier
}

% Comando para insertar una imagen sin posición.
% Los parámetros son:
% 1 --> Ruta absoluta/relativa de la figura
% 2 --> Texto a pie de figura
% 3 --> Tamaño en tanto por uno relativo al ancho de página
\newcommand{\imagenflotante}[3]{
	\begin{figure}
		\centering
		\includegraphics[width=#3\textwidth]{#1}
		\caption{#2}\label{fig:#1}
	\end{figure}
}

% El comando \figura nos permite insertar figuras comodamente, y utilizando
% siempre el mismo formato. Los parametros son:
% 1 --> Porcentaje del ancho de página que ocupará la figura (de 0 a 1)
% 2 --> Fichero de la imagen
% 3 --> Texto a pie de imagen
% 4 --> Etiqueta (label) para referencias
% 5 --> Opciones que queramos pasarle al \includegraphics
% 6 --> Opciones de posicionamiento a pasarle a \begin{figure}
\newcommand{\figuraConPosicion}[6]{%
  \setlength{\anchoFloat}{#1\textwidth}%
  \addtolength{\anchoFloat}{-4\fboxsep}%
  \setlength{\anchoFigura}{\anchoFloat}%
  \begin{figure}[#6]
    \begin{center}%
      \Ovalbox{%
        \begin{minipage}{\anchoFloat}%
          \begin{center}%
            \includegraphics[width=\anchoFigura,#5]{#2}%
            \caption{#3}%
            \label{#4}%
          \end{center}%
        \end{minipage}
      }%
    \end{center}%
  \end{figure}%
}

%
% Comando para incluir imágenes en formato apaisado (sin marco).
\newcommand{\figuraApaisadaSinMarco}[5]{%
  \begin{figure}%
    \begin{center}%
    \includegraphics[angle=90,height=#1\textheight,#5]{#2}%
    \caption{#3}%
    \label{#4}%
    \end{center}%
  \end{figure}%
}
% Para las tablas
\newcommand{\otoprule}{\midrule [\heavyrulewidth]}
%
% Nuevo comando para tablas pequeñas (menos de una página).
\newcommand{\tablaSmall}[5]{%
 \begin{table}
  \begin{center}
   \rowcolors {2}{gray!35}{}
   \begin{tabular}{#2}
    \toprule
    #4
    \otoprule
    #5
    \bottomrule
   \end{tabular}
   \caption{#1}
   \label{tabla:#3}
  \end{center}
 \end{table}
}

%
% Nuevo comando para tablas pequeñas (menos de una página).
\newcommand{\tablaSmallSinColores}[5]{%
 \begin{table}[H]
  \begin{center}
   \begin{tabular}{#2}
    \toprule
    #4
    \otoprule
    #5
    \bottomrule
   \end{tabular}
   \caption{#1}
   \label{tabla:#3}
  \end{center}
 \end{table}
}

\newcommand{\tablaApaisadaSmall}[5]{%
\begin{landscape}
  \begin{table}
   \begin{center}
    \rowcolors {2}{gray!35}{}
    \begin{tabular}{#2}
     \toprule
     #4
     \otoprule
     #5
     \bottomrule
    \end{tabular}
    \caption{#1}
    \label{tabla:#3}
   \end{center}
  \end{table}
\end{landscape}
}

%
% Nuevo comando para tablas grandes con cabecera y filas alternas coloreadas en gris.
\newcommand{\tabla}[6]{%
  \begin{center}
    \tablefirsthead{
      \toprule
      #5
      \otoprule
    }
    \tablehead{
      \multicolumn{#3}{l}{\small\sl continúa desde la página anterior}\\
      \toprule
      #5
      \otoprule
    }
    \tabletail{
      \hline
      \multicolumn{#3}{r}{\small\sl continúa en la página siguiente}\\
    }
    \tablelasttail{
      \hline
    }
    \bottomcaption{#1}
    \rowcolors {2}{gray!35}{}
    \begin{xtabular}{#2}
      #6
      \bottomrule
    \end{xtabular}
    \label{tabla:#4}
  \end{center}
}

%
% Nuevo comando para tablas grandes con cabecera.
\newcommand{\tablaSinColores}[6]{%
  \begin{center}
    \tablefirsthead{
      \toprule
      #5
      \otoprule
    }
    \tablehead{
      \multicolumn{#3}{l}{\small\sl continúa desde la página anterior}\\
      \toprule
      #5
      \otoprule
    }
    \tabletail{
      \hline
      \multicolumn{#3}{r}{\small\sl continúa en la página siguiente}\\
    }
    \tablelasttail{
      \hline
    }
    \bottomcaption{#1}
    \begin{xtabular}{#2}
      #6
      \bottomrule
    \end{xtabular}
    \label{tabla:#4}
  \end{center}
}

%
% Nuevo comando para tablas grandes sin cabecera.
\newcommand{\tablaSinCabecera}[5]{%
  \begin{center}
    \tablefirsthead{
      \toprule
    }
    \tablehead{
      \multicolumn{#3}{l}{\small\sl continúa desde la página anterior}\\
      \hline
    }
    \tabletail{
      \hline
      \multicolumn{#3}{r}{\small\sl continúa en la página siguiente}\\
    }
    \tablelasttail{
      \hline
    }
    \bottomcaption{#1}
  \begin{xtabular}{#2}
    #5
   \bottomrule
  \end{xtabular}
  \label{tabla:#4}
  \end{center}
}



\definecolor{cgoLight}{HTML}{EEEEEE}
\definecolor{cgoExtralight}{HTML}{FFFFFF}

%
% Nuevo comando para tablas grandes sin cabecera.
\newcommand{\tablaSinCabeceraConBandas}[5]{%
  \begin{center}
    \tablefirsthead{
      \toprule
    }
    \tablehead{
      \multicolumn{#3}{l}{\small\sl continúa desde la página anterior}\\
      \hline
    }
    \tabletail{
      \hline
      \multicolumn{#3}{r}{\small\sl continúa en la página siguiente}\\
    }
    \tablelasttail{
      \hline
    }
    \bottomcaption{#1}
    \rowcolors[]{1}{cgoExtralight}{cgoLight}

  \begin{xtabular}{#2}
    #5
   \bottomrule
  \end{xtabular}
  \label{tabla:#4}
  \end{center}
}



\graphicspath{ {./img/} }

% Capítulos
\chapterstyle{bianchi}
\newcommand{\capitulo}[2]{
	\setcounter{chapter}{#1}
	\setcounter{section}{0}
	\setcounter{figure}{0}
	\setcounter{table}{0}
	\chapter*{\thechapter.\enskip #2}
	\addcontentsline{toc}{chapter}{\thechapter.\enskip #2}
	\markboth{#2}{#2}
}

% Apéndices
\renewcommand{\appendixname}{Apéndice}
\renewcommand*\cftappendixname{\appendixname}

\newcommand{\apendice}[1]{
	%\renewcommand{\thechapter}{A}
	\chapter{#1}
}

\renewcommand*\cftappendixname{\appendixname\ }

% Formato de portada
\makeatletter
\usepackage{xcolor}
\newcommand{\tutor}[1]{\def\@tutor{#1}}
\newcommand{\course}[1]{\def\@course{#1}}
\definecolor{cpardoBox}{HTML}{E6E6FF}
\def\maketitle{
  \null
  \thispagestyle{empty}
  % Cabecera ----------------
\noindent\includegraphics[width=\textwidth]{cabecera}\vspace{1cm}%
  \vfill
  % Título proyecto y escudo informática ----------------
  \colorbox{cpardoBox}{%
    \begin{minipage}{.8\textwidth}
      \vspace{.5cm}\Large
      \begin{center}
      \textbf{TFG del Grado en Ingeniería Informática}\vspace{.6cm}\\
      \textbf{\LARGE\@title{}}
      \end{center}
      \vspace{.2cm}
    \end{minipage}

  }%
  \hfill\begin{minipage}{.20\textwidth}
    \includegraphics[width=\textwidth]{escudoInfor}
  \end{minipage}
  \vfill
  % Datos de alumno, curso y tutores ------------------
  \begin{center}%
  {%
    \noindent\LARGE
    Presentado por \@author{}\\ 
    en Universidad de Burgos --- \@date{}\\
    Tutor: \@tutor{}\\
  }%
  \end{center}%
  \null
  \cleardoublepage
  }
\makeatother

\newcommand{\nombre}{César Rodríguez Villagrá} %%% cambio de comando

% Datos de portada
\title{GII-24.19-contramedidas-IoT-mediante-reinforcement-learning}
\author{\nombre}
\tutor{nombre tutor}
\date{\today}

\begin{document}

\maketitle


\newpage\null\thispagestyle{empty}\newpage


%%%%%%%%%%%%%%%%%%%%%%%%%%%%%%%%%%%%%%%%%%%%%%%%%%%%%%%%%%%%%%%%%%%%%%%%%%%%%%%%%%%%%%%%
\thispagestyle{empty}


\noindent\includegraphics[width=\textwidth]{cabecera}\vspace{1cm}

\noindent D. nombre tutor, profesor del departamento de nombre departamento, área de nombre área.

\noindent Expone:

\noindent Que el alumno D. \nombre, con DNI dni, ha realizado el Trabajo final de Grado en Ingeniería Informática titulado título de TFG. 

\noindent Y que dicho trabajo ha sido realizado por el alumno bajo la dirección del que suscribe, en virtud de lo cual se autoriza su presentación y defensa.

\begin{center} %\large
En Burgos, {\large \today}
\end{center}

\vfill\vfill\vfill

% Author and supervisor
\begin{minipage}{0.45\textwidth}
\begin{flushleft} %\large
Vº. Bº. del Tutor:\\[2cm]
D. nombre tutor
\end{flushleft}
\end{minipage}
\hfill
\begin{minipage}{0.45\textwidth}
\begin{flushleft} %\large
Vº. Bº. del co-tutor:\\[2cm]
D. nombre co-tutor
\end{flushleft}
\end{minipage}
\hfill

\vfill

% para casos con solo un tutor comentar lo anterior
% y descomentar lo siguiente
%Vº. Bº. del Tutor:\\[2cm]
%D. nombre tutor


\newpage\null\thispagestyle{empty}\newpage




\frontmatter

% Abstract en castellano
\renewcommand*\abstractname{Resumen}
\begin{abstract}
En este primer apartado se hace una \textbf{breve} presentación del tema que se aborda en el proyecto.
\end{abstract}

\renewcommand*\abstractname{Descriptores}
\begin{abstract}
Palabras separadas por comas que identifiquen el contenido del proyecto Ej: servidor web, buscador de vuelos, android \ldots
\end{abstract}

\clearpage

% Abstract en inglés
\renewcommand*\abstractname{Abstract}
\begin{abstract}
A \textbf{brief} presentation of the topic addressed in the project.
\end{abstract}

\renewcommand*\abstractname{Keywords}
\begin{abstract}
keywords separated by commas.
\end{abstract}

\clearpage

% Indices
\tableofcontents

\clearpage

\listoffigures

\clearpage

\listoftables
\clearpage

\mainmatter
\capitulo{1}{Introducción}

Descripción del contenido del trabajo y del estructura de la memoria y del resto de materiales entregados.

\capitulo{2}{Objetivos del proyecto}
El objetivo general del proyecto es el desarrollo de un agente de aprendizaje por refuerzo capaz de mitigar ataques de denegación de servicio basados en la ocupación de la cola de un router de la red, en un entorno de dispositivos de internet de las cosas. Para ello, se ha desarrollado un entorno de simulación que permite la generación de tráfico benigno y malicioso, así como el entrenamiento del agente. Además, se ha desarrollado una aplicación web que permite la gestión de los modelos entrenados por cada usuario, mostrar la documentación asociada y visualizar los resultados obtenidos por el agente, junto con la visualización de la función de recompensa y el estado del entorno.

Los principales objetivos del proyecto son:
\begin{itemize}
    \item Implementar un entorno de simulación para el entrenamiento del agente.
    \item Desarrollar un agente de inteligencia artificial, mediante aprendizaje por refuerzo capaz de detectar y mitigar ataques a una red informática, especialmente en un entorno IoT.
    \item Desarrollar y desplegar una aplicación web que permita la visualización de los resultados obtenidos por el agente.
\end{itemize}

\section{Objetivos específicos}
Los objetivos más específicos del proyecto son:

\subsection{Generación del entorno de simulación y la generación de los datos de entrenamiento}
El objetivo de esta parte del proyecto es la generación artificial de generadores de paquetes para simular tráfico benigno y tráfico atacante en un entorno de dispositivos de internet de las cosas, así como la creación de un entorno de simulación para el entrenamiento del agente de aprendizaje por refuerzo. Los requisitos específicos de este apartado son:
\begin{itemize}
    \item Desarrollar generadores de tráfico benigno representativo de dispositivos IoT basado en datos realistas.
    \item Implementar generadores de tráfico malicioso que simulen ataques de denegación de servicio basado en datos realistas.
    \item Realizar un estudio de la función de recompensa para su optimización, en el que se muestre la evolución de la recompensa según la acción y el estado del entorno.
    \item Crear un entorno de simulación que permita el entrenamiento y la evaluación del agente.
    \item Realizar una evaluación gráfica de los resultados obtenidos por el agente.
\end{itemize}

\subsection{Desarrollo de la aplicación web}
El propósito de este apartado es el desarrollo de una aplicación web para la gestión y visualización de los modelos y resultados. Los requisitos específicos son:
\begin{itemize}
    \item Implementar una interfaz para la visualización de algunos los resultados obtenidos por el agente.
    \item Permitir la gestión y almacenamiento persistente de los modelos entrenados por los usuarios.
    \item Incluir la documentación asociada a cada modelo, accesible desde la aplicación.
    \item Facilitar la visualización de la función de recompensa y del estado del entorno de simulación.
    \item Permitir la gestión de los modelos almacenados por cada usuario.
    \item Permitir la gestión de los usuarios por parte del administrador.
    \item Proporcionar ayuda a los usuarios en la aplicación web para facilitar su uso.
\end{itemize}

\subsection{Otros objetivos}
Además de los objetivos ya descritos, es importante destacar otros objetivos que se han tenido en cuenta durante el desarrollo del proyecto:
\begin{itemize}
    \item \textbf{Seguridad:} Desarrollar el sistema siguiendo buenas prácticas de seguridad informática para proteger los datos y la integridad del sistema.
    \item \textbf{Reducción de costes:} Diseñar el sistema de forma que se minimicen los recursos requeridos y el costo asociado, especialmente para el despliegue de la aplicación web.
    \item \textbf{Usabilidad:} Diseñar una interfaz web intuitiva, accesible y simple para los usuarios, asegurando una experiencia de usuario satisfactoria.
    \item \textbf{Escalabilidad:} Permitir que el sistema pueda ser fácilmente ampliado o adaptado a nuevas necesidades o requisitos futuros.
    \item \textbf{Mantenibilidad:} Estructurar el código del proyecto de forma modular y documentada para facilitar su mantenimiento, extensión o reutilización en futuros desarrollos, siguiendo patrones de diseño y buenas prácticas de programación.
\end{itemize}
\capitulo{3}{Conceptos teóricos}
En este capítulo se explican los conceptos teóricos relacionados con el proyecto, necesarios para la comprensión del mismo.

\section{Aprendizaje por refuerzo}
El aprendizaje por refuerzo o \textit{Reinforcement Learning} (RL) es una rama del aprendizaje automático que su objetivo es aprender cómo los agentes deben tomar decisiones en un entorno, para maximizar la recompensa asociada a sus acciones y estado del entorno. A diferencia del aprendizaje supervisado, donde se dispone de un conjunto de datos etiquetados, en el aprendizaje por refuerzo el agente aprende a través de la interacción con el entorno, recibiendo recompensas o penalizaciones en función de las acciones que realiza. Su aprendizaje se basa en la prueba y error, donde el agente aprende a través de la experiencia acumulada en el entorno.
Los principales elementos del aprendizaje por refuerzo son:
\begin{itemize}
	\item \textbf{Agente:} Es el ente que se encarga de la toma de decisiones y aprende a través de la interacción con el entorno a lo largo del tiempo.
	\item \textbf{Entorno:} Es el contexto en el que el agente se encuentra, incluyendo todos los elementos con los que se relaciona.
	\item \textbf{Estado:} Es una instantánea en un momento dado, que puede incluir información sobre el estado actual del agente y del entorno.
	\item \textbf{Acción:} Es una decisión tomada por el agente que afecta al estado del entorno.
	\item \textbf{Política:} Es una estrategia que define cómo el agente selecciona acciones en función del estado actual del entorno.
	\item \textbf{Recompensa:} Es una señal que indica al agente la calidad de la acción tomada en un estado determinado. El objetivo del agente es maximizar la recompensa acumulada a lo largo del tiempo.
\end{itemize}

El aprendizaje por refuerzo se utiliza en una amplia variedad de aplicaciones, como juegos, robótica, optimización de procesos, entre otros.

Se usa en situaciones donde:
\begin{itemize}
	\item Las acciones afectan el estado futuro del entorno
	\item No se dispone de un conjunto de datos etiquetados
	\item Se tiene retroalimentación escasa o tardía de la acción que se ha tomado
	\item El agente debe aprender a través de la experiencia acumulada.
\end{itemize}
Es especialmente útil en entornos muy variables o dinámicos, donde se tienen que tomar acciones secuencialmente y se necesita una adaptación a los cambios en el entorno.


\section{Redes de computadoras}
Una red de computadoras es un conjunto de dispositivos interconectados que pueden comunicarse entre sí. Permiten el intercambio de mensajes y recursos entre los elementos de la red. Actualmente las redes de computadoras están muy extendidas y son fundamentales para el funcionamiento de Internet y de muchas aplicaciones y servicios en línea.
Las redes se pueden clasificar en diferentes tipos según su alcance, los más comunes son:
\begin{itemize}
	\item \textbf{Redes de área local (LAN):} Conectan computadoras en un área geográfica limitada, como una oficina o un edificio.
	\item \textbf{Redes de área amplia (WAN):} Conectan computadoras en áreas geográficas más amplias, como ciudades o países.
\end{itemize}

Los elementos principales de una red de computadoras son:
\begin{itemize}
	\item \textbf{Dispositivos finales:} Son los dispositivos que se conectan a la red, los cuales se quieren comunicar con otros dispositivos finales. Los dispositivos finales pueden ser computadoras, teléfonos móviles, impresoras, cámaras, sensores, etc.
	\item \textbf{Dispositivos de red:} Son los dispositivos que permiten la comunicación entre los dispositivos finales, como routers, switches, etc.
	\item \textbf{Medios de transmisión:} Son los canales a través de los cuales se transmiten la información entre dispositivos, como cables Ethernet, fibra óptica, Wi-Fi, etc.
\end{itemize}

La comunicación entre dispositivos se realiza a través de paquetes de datos que viajan por los medios de transmisión. Los paquetes incluyen información sobre el origen, destino, protocolo, y contenido del mensaje entre otros. Es importante tener en cuenta el tamaño de los paquetes, siendo necesario en algunas ocasiones dividir los mensajes en varios paquetes para su transmisión.
Los protocoles de red son un pilar fundamental en la comunicación de computadoras en redes, ya que son un conjunto de reglas y convenciones que definen cómo se transmiten los datos este estos.
Algunos de protocolos más comunes son:
\begin{itemize}
	\item \textbf{HTTP (HyperText Transfer Protocol):} Es el protocolo que se usa par la transferencia de datos en internet. Permite la comunicación entre navegadores web y servidores. Sin embargo, HTTP tiene problemas de seguridad, ya que los datos no se transmiten cifrados.
	\item \textbf{TLS (Transport Layer Security):} Es un protocolo criptográfico que proporciona comunicaciones seguras sobre una red, su predecesesor es la capa de puertos seguros o SSL. TLS cifra los datos antes de enviarlos y verifica que no hayan sido alterados durante la transmisión. Es ampliamente utilizado en protocolos como HTTPS, correo electrónico, y otros servicios en línea. ~\cite{SushilJajodia2025}

	\item \textbf{HTTPS (HTTP Secure):} Es la versión de tranmisión de datos segura de HTTP. Utiliza una capa adicional de seguridad llamada TLS (Transport Layer Security) para cifrar la información durante la transmisión. Esto garantiza la seguridad de los datos, protegiendo contra ataques como la interceptación o modificación de mensajes. ~\cite{SushilJajodia2025}

	\item \textbf{TCP/IP (Transmission Control Protocol / Internet Protocol):} Es el conjunto de protocolos más utilizado en Internet. TCP se encarga de dividir los mensajes en paquetes, enviarlos y asegurarse de que lleguen completos y en orden. El protocolo TCP está hecho para ser utilizado con el IP, que se encarga de direccionar y enrutar los paquetes hacia su destino a través de la red.~\cite{Protocolo}

	\item \textbf{IPv4 (Internet Protocol version 4) y IPv6 (Internet Protocol version 6):} Es la cuarta versión del protocolo Internet Protocol(IP) y la más utilizada. Utiliza direcciones de 32 bits. Debido al crecimiento de dispositivos conectados, se ha ido agotando el espacio de direcciones IPv4.
	IPv6 Es la versión más reciente del protocolo IP. Utiliza direcciones de 128 bits, lo que soluciona el problema de las direcciones limitadas IPv4. IPv6 también mejora aspectos como la eficiencia y configuración entre otros.

	\item \textbf{DNS (Domain Name System):} Es el sistema que traduce los nombres de dominio en direcciones IP, facilitando al usuario el acceso a los servidores que proporcionan la información solicitada. Principalmente permite que los usuarios accedan a sitios web usando nombres fáciles de recordar(dominios) en lugar de números IP difíciles de manejar.

\end{itemize}


En cuanto a los dispositivos con funciones de seguridad de red, los principales son:
\begin{itemize}
	\item \textbf{Firewalls:} Son dispositivos hardware o software que controlan el tráfico de red, permitiendo o bloqueando el paso de paquetes según reglas predefinidas. Se utilizan para proteger la red de accesos no autorizados y ataques informáticos. ~\cite{SushilJajodia2025}

	\item \textbf{Proxies:} Como Cloudflare, son servidores que actúan como intermediarios entre los dispositivos finales y los servidores de destino. Se utilizan para mejorar el rendimiento, la seguridad y la privacidad de la red. ~\cite{SushilJajodia2025}

	\item \textbf{CDN (Content Delivery Network):} Son redes de servidores distribuidos en distintas ubicaciones que almacenan copias del contenido estático (como imágenes, scripts etc) para entregarlo de forma más rápida y eficiente al usuario. Además de mejorar el rendimiento, muchas CDN también contribuyen a la seguridad al absorber tráfico malicioso. Cloudflare es un ejemplo de CDN que ofrece servicios de seguridad y rendimiento para aplicaciones web.
\end{itemize}

\section{Interfaces de Programación de Aplicaciones}
Una API (\textit{Application Programming Interface}) es un conjunto de definiciones y protocolos que permiten que dos aplicaciones se comuniquen entre sí. Las APIs actúan como una interfaz que facilita el acceso a funcionalidades o datos de un sistema sin necesidad de conocer su implementación interna.~\cite{Api}
Una API REST (\textit{Representational State Transfer}) es un tipo de API que sigue los principios de la arquitectura REST, siguiendo estos principios:
\begin{itemize}
	\item \textbf{Cliente-servidor:} La API REST separa el cliente del servidor, permitiendo que ambos se desarrollen de forma independiente, pero manteniendo la comunicación entre ellos.
	\item \textbf{Sin estado:} Donde en el servidor no se almacenen datos del cliente entre solicitudes.
	\item \textbf{Cacheable:} Las respuestas de la API REST deben poderse guardar en caché, para mejorar el rendimiento y reducir la carga en el servidor.
	\item \textbf{Interfaz uniforme:} La API REST debe tener una interfaz consistente y fácil de entender, para facilitar su uso y mantenimiento.
\end{itemize}
Estas llamadas a la API REST se realizan a través de peticiones HTTP, utilizando los métodos más comunes del protocolo, como GET para obtener datos del servidor, POST para enviar o crear datos, PUT para actualizar recursos existentes, y DELETE para eliminar recursos. Existen otros métodos menos utilizados, pero estos son los más frecuentes.

\section{Internet de las cosas}
Los dispositivos de Internet de las cosas, o por sus siglas en inglés \textit{Internet of Things} (IoT), son dispositivos físicos que están conectados a una red y se encargan de enviar, recibir y procesar datos del mundo real. ~\cite{AzizAlKabir2023} 

Actualmente estos dispositivos están muy extendidos y son fundamentales para el funcionamiento de muchas aplicaciones y servicios en línea, como la domótica, la monitorización de la salud, la gestión de la energía, cámaras de seguridad, etc. Estos dispositivos se caracterizan por la escasa capacidad de procesamiento y almacenamiento, siendo también muy susceptibles a ataques informáticos debido a su baja seguridad. Un ataque muy común en estos dispositivos es el ataque de denegación de servicio, que busca saturar los recursos del dispositivo o de la red a la que está conectado, impidiendo su correcto funcionamiento, o incluso utilizar estos dispositivos para formar botnets. ~\cite{kolias2017}


\section{Ataques informáticos}
Los ataques informáticos son acciones que tienen como objetivo comprometer la seguridad de los sistemas, redes o datos.
Estas acciones, su objetivo es afectar a uno o varios de los tres pilares fundamentales de la ciberseguridad: la confidencialidad, integridad y disponibilidad (conocidos como el triángulo CIA).~\cite{NicholasEdwards2020, Chaeikar2012} 

Actualmente los ataques informáticos son una amenaza constante para la seguridad de las redes y sistemas. Estos ataques pueden tener diferentes objetivos, como robar información, interrumpir servicios, entre otros.

\begin{itemize}
    \item \textbf{Ataques de denegación de servicio (DoS y DDoS):} Buscan saturar los recursos de un sistema o red, impidiendo el acceso de los usuarios legítimos a los servicios. Denied of Service (DoS) se refiere a un ataque desde una sola fuente, mientras que Distributed Denied of Service (DDoS) implica la versión distribuida, atacando desde múltiples fuentes. Para ataques DDoS se suelen usar redes de bots o \textit{botnets}, que son redes de dispositivos comprometidos que se controlan de forma remota para llevar a cabo el ataque. Debido a la gran cantidad de dispositivos de internet de las cosas que se usan actualmente, y su baja seguridad, son un objetivo común para los atacantes que buscan crear redes de bots. ~\cite{kolias2017}

    \item \textbf{Malware:} Es un programa informático malicioso diseñado para dañar, robar o modificar la información o sistemas informáticos. Incluye virus, gusanos, troyanos, ransomware, entre otros. ~\cite{SushilJajodia2025}

    \item \textbf{Exploits de vulnerabilidades:} Son técnicas o herramientas que aprovechan de fallos de seguridad en el software o hardware para comprometer el sistema. ~\cite{SushilJajodia2025}

	\item \textbf{Ataques CSRF(Cross-Site Request Forgery):} CSRF es un tipo de ataque que fuerza a un usuario autenticado en el sistema para que realice acciones no deseadas. El atacante envía una solicitud al navegador del usuario, que se ejecuta sin su conocimiento, aprovechando la sesión activa del usuario.

	\item \textbf{Ataques XSS(Cross-Site Scripting):} XSS es un tipo de ataque que permite a un atacante incluir scripts maliciosos en una página web, que se ejecutan en el navegador de los usuarios que visitan esa página. Estos scripts pueden robar información del usuario, o realizar acciones no autorizadas en nombre del usuario.
\end{itemize}

Para que los ataques informáticos no tengan éxito, es importante implementar medidas de seguridad adecuadas y mitigar los ataques que se detecten en la medida de lo posible. 
Algunas de las mitigaciones más comunes son:
\begin{itemize}
	\item \textbf{Descartar paquetes:} Se descartan los paquetes que se consideran basados en reglas establecidas, por ejemplo, paquetes que no cumplen con el formato esperado, que provienen de ciertas direcciones IP, ubicaciones geográficas, protocolos etc. ~\cite{Yungaicela-Naula2022, Liu2018} 

    \item \textbf{Rate limiting:} Consiste en limitar la cantidad de peticiones que puede realizar un cliente en un intervalo de tiempo determinado. Esto permite evitar que el servidor se pueda saturar por un gran número de peticiones desde un mismo cliente, ya sea un cliente legítimo o un bot malicioso que realice escaneos masivos, ataques de denegación de servicio o ataques de fuerza bruta.

    \item \textbf{Desafíos interactivos (tipo CAPTCHA):} Es común encontrar en aplicaciones web desafíos para verificar si el usuario que intenta acceder al recurso es humano o no. Estos desafíos suelen ser imágenes, preguntas o tareas que requieren la intervención del usuario para conseguir pasarlos.~\cite{CloudflareDocsTeam2025}
\end{itemize}
\capitulo{4}{Técnicas y herramientas}
Las técnicas y herramientas utilizadas en el desarrollo de un proyecto son fundamentales para cumplir su objetivo de manera eficiente. En este apartado se describen las principales tecnologías, lenguajes de programación, bibliotecas y herramientas de desarrollo empleadas en el proyecto, así como su justificación y el papel que desempeñan en el mismo.

Los principales lenguajes de programación utilizados en el proyecto son Python y Java.

Python se ha utilizado principalmente para el desarrollo del agente de aprendizaje por refuerzo, debido a su amplia gama de bibliotecas y frameworks especializados en este campo, como Stable-Baselines3~\cite{Raffin2021} y Gymnasium, también se ha usado para el desarrollo de la api con Flask-RESTful, para facilitar la posible interacción de la web con el agente, y que permite una implementación rápida y eficiente en recursos.

Java se ha utilizado para el desarrollo de la aplicación web, aprovechando su robustez y escalabilidad, junto con el framework Spring Boot 3 para facilitar la creación de aplicaciones web, para la gestión de los paquetes asociados se ha utilizado Maven, que permite una gestión eficiente de las dependencias y la construcción del proyecto. Para las pruebas unitarias se ha utilizado JUnit, y para los informes de pruebas Jacoco. Spring Boot tiene una gran comunidad y soporte, lo que facilita la resolución de problemas y la implementación de nuevas funcionalidades, junto con la posibilidad de integrar fácilmente otras tecnologías y herramientas, como es el caso de persistencia de datos con JPA, Hibernate y MySQL, la integración de un sistema de autenticación y autorización robusto con Spring Security, o la integración de un modelo de lenguaje a través de una API.

Los principales entornos de desarrollo utilizados han sido:
\begin{itemize}
    \item \textbf{IntelliJ IDEA:} Para el desarrollo de la aplicación web en Java, que ofrece una amplia gama de herramientas y características para facilitar el desarrollo, refactorización, depuración, entre otras. S ha utilizado con la extensión de SonarQube.
    \item \textbf{Visual Studio Code:} Para el desarrollo del agente de aprendizaje por refuerzo en Python y la API REST, que es un editor ligero y altamente configurable. Se ha utilizado con extensiones para Python, Jupyter Notebook, SonarQube y Docker.
\end{itemize}

Para la visualización de datos en la aplicación web se ha utilizado Plotly.js, que permite crear gráficos interactivos y visualizaciones avanzadas de datos, mejorando la experiencia del usuario. Para el desarrollo del frontend se ha utilizado HTML5, CSS3 y JavaScript, junto con Bootstrap 5 para facilitar el diseño responsivo y moderno de la interfaz de usuario.

Como control de versiones se ha usado Git, que permite llevar un seguimiento de los cambios realizados en el código fuente del proyecto, facilitando la gestión de versiones, en el que el cliente usado es GitHub, permitiendo el almacenamiento de las versiones en la nube, junto con la posibilidad de usar GitHub Actions, e integrarlo con GitHub Projects para la gestión de tareas y seguimiento del proyecto. 

Para la gestión de la documentación se ha utilizado LaTeX, que permite crear documentos técnicos de alta calidad, y para la transformación a pdf se ha utilizado una GitHub Action que usa una imagen de Docker para transformar los archivos fuente a pdf, y para la gestión de referencias bibliográficas se ha utilizado Mendeley, que facilita la gestión y citación de referencias en el proyecto.

Para la creación de diagramas y gráficos se ha utilizado Draw.io, que permite crear todos los diagramas de manera sencilla y rápida, con integración tanto en la web como en Visual Studio Code.

Para la calidad de código y la detección de errores se ha utilizado SonarQube, que permite analizar el código fuente del proyecto y detectar posibles errores, vulnerabilidades y problemas de calidad. Se ha integrado con GitHub Actions para realizar análisis automáticos del código en cada commit y pull request. también se han utilizado AutoPep8 para el formateo del código Python, que permiten mantener un estilo de codificación consistente y mejorar la legibilidad del código.

Para la gestión de los registros DNS, compra del dominio y seguridad de la aplicación web se ha utilizado Cloudflare. Para la compra del dominio se ha elegido Cloudflare debido a que ofrece un servicio de registro de dominios a precios competitivos y con una interfaz fácil de usar, además de que proporciona una gran varidad de servicios, reduciendo la necesidad de utilizar múltiples proveedores para diferentes servicios. Cloudflare proporciona en su capa gratuita amplios servicios de seguridad y rendimiento para aplicaciones web, incluyendo protección contra ataques DDoS, optimización del rendimiento y gestión de certificados SSL. Cloudflare es una plataforma muy popular y confiable, utilizada por muchas empresas, particulares y organizaciones en todo el mundo, con una alta efectividad. ~\cite{Nadeem2023,Adhar2023}

En el despliegue de la aplicación se ha elegido un servidor AWS EC2, que permite un control total sobre la infraestructura y la configuración del servidor. Es una de las principales plataformas de computación en la nube más populares, confiables y seguras, con una amplia gama de servicios y características que permiten escalar y gestionar aplicaciones de manera eficiente. La capa gratuita permite la ejecución de aplicaciones pequeñas pero con funcionalidades limitadas. La capa gratuita, proporciona una instancia con 1GB de RAM, lo cual no es susficiente para ejecutar este proyecto, por lo que se ha optado por una isntancia t3.medium, que proporciona 4GB de RAM a un precio asequible.

En la siguiente tabla se muestran marcadas con X las herramientas y tecnologías utilizadas en cada parte del proyecto. \tabla{Herramientas y tecnologías utilizadas en cada parte del proyecto}{l c c c c c c}{7}{herramientasportipodeuso}
{ \multicolumn{1}{l}{Herramientas} & \shortstack{Modelo de\\aprendizaje\\por refuerzo}  & Web & API REST & BD & Memoria &  \\}{ 
Python                 & X &   & X &   &   \\
Jupyter Notebook       & X &   &   &   &   \\
Stable-Baselines3      & X &   &   &   &   \\
Gymnasium              & X &   &   &   &   \\
Matplotlib             & X &   &   &   &   \\
Java                   &   & X &   &   &   \\
Spring Boot 3          &   & X &   &   &   \\
Maven                  &   & X &   &   &   \\
JUnit                  &   & X &   &   &   \\
Selenium               &   & X &   &   &   \\
HTML5                  &   & X &   &   &   \\
CSS3                   &   & X &   &   &   \\
Plotly.js              &   & X &   &   &   \\
JavaScript             &   & X &   &   &   \\
Bootstrap 5            &   & X &   &   &   \\
Flask-RESTful          &   &   & X &   &   \\
Pytest                 &   &   & X &   &   \\
REST                   &   & X & X &   &   \\
JSON                   & X & X & X &   &   \\
Docker                 & X & X & X & X &   \\
Docker Compose         & X & X & X & X &   \\
MySQL                  &   &   &   & X &   \\
SQL                    & X &   &   & X &   \\
Git                    & X & X & X & X & X \\
GitHub Projects        & X & X & X & X & X \\
GitHub Actions         & X & X & X &   & X \\
GitHub Dependabot      & X & X & X &   & X \\
GitHub Secret Scanning & X & X & X &   & X \\
GitHub Code Scanning   & X & X & X &   & X \\
GitHub Copilot         & X & X & X &   & X \\
SonarQube              & X & X & X &   & X \\
Draw.io                &   &   &   &   & X \\
LaTeX                  &   &   &   &   & X \\
IntelliJ IDEA          &   & X &   & X &   \\
VS Code                & X &   & X &   & X \\
Mendeley               &   &   &   &   & X \\
AWS EC2                &   & X & X & X &   \\
Cloudflare             &   & X &   &   &   \\
API Groq               &   & X &   &   &   \\
} 
\capitulo{5}{Aspectos relevantes del desarrollo del proyecto}

El proyecto se ha dividido en tres partes, cada una de las cuales tiene un objetivo específico:
\begin{itemize}
    \item \textbf{Generación del entorno de simulación y la generación de los datos de entrenamiento:} Cuyo objetivo es la generación artificial de paquetes para simular tráfico atacante y tráfico de benigno en un entorno de dispositivos de internet de las cosas. También tiene como objetivo la creación de un entorno de simulación para el entrenamiento del agente de aprendizaje por refuerzo.
    \item \textbf{Desarrollo del entorno y el agente de aprendizaje por refuerzo:} Incluyendo el ajuste de parámetros, mejora en la función de recompensa, entrenamiento del agente y su posterior evaluación una vez entrenado.
    \item \textbf{Desarrollo de la aplicación web:} Su principal objetivo es la visualización de los resultados obtenidos por el agente, y que los usuarios puedan visualizar de manera sencilla los resultados.
\end{itemize}

\section{Generación del entorno de simulación y la generación de los datos de entrenamiento}
Para el entorno de simulación se ha utilizado la libería de Python \texttt{Gymnasium}, que permite facilita la creación del entorno de simulación para el agente de aprendizaje por refuerzo.
La arquitectura del entorno es la siguiente:
\imagen{RedInicial}{Arquitectura del entorno de simulación}{1.0}

La generación de los datos de entrenamiento se ha realizado creando generadores de paquetes, tanto de tráfico benigno como de tráfico atacante, que simulan el tráfico que se puede generar en un entorno de dispositivos de internet de las cosas. Los atributos de los paquetes se generan siguiendo una distribución por probabilidades, para que se asemeje lo más posible al tráfico real que se puede generar en un entorno de dispositivos de internet de las cosas, tanto, cuando reciben un ataque de denegación de servicio, como cuando no.

En el entorno de entrenamiento los estados de ataque y benigno cambian siguiendo una máquina de estados, en la que hay una cierta probabilidad de que el estado cambie. En cuanto entre cada cambio de estado los paquetes transmitidos tienen las mismas características, es decir cada ataque tiene un patrón de tráfico específico que cambia cuando cambia de estado. Esto permite que los ataques sean variables para que el agente pueda aprender a detectar ataques de denegación de servicio con diferentes características.

\section{Desarrollo del entorno y el agente de aprendizaje por refuerzo}

\subsection{Algoritmo de entrenamiento del agente}
En el entrenamiento del agente se ha usado el algoritmo de aprendizaje por refuerzo Proximal Policy Optimization (PPO), que es un algoritmo de optimización de políticas que se basa en la idea de maximizar la función de recompensa esperada, y que ha demostrado ser efectivo en una amplia variedad de entornos y tareas.~\cite{Siboo2023, Yu2022}
El algoritmo PPO se encuentra implementado en la librería de Python \texttt{Stable Baselines3}, que es una librería de aprendizaje por refuerzo que proporciona implementaciones de algoritmos de aprendizaje por refuerzo basados en PyTorch, y que es ampliamente utilizada en la comunidad de aprendizaje por refuerzo.~\cite{Raffin2021}
Este tiene una compatibilidad total con la librería \texttt{Gymnasium}, que es la librería que se ha usado para crear el entorno de simulación, y con distintos tipos de espacios, que son cómo se han definido los estados, las acciones y las recompensas del entorno de simulación. En el caso de PPO tiene compatibilidad con espacios todos los espacios de \texttt{Gymnasium}, incluyendo espacios Box, Discrete, MultiDiscrete y MultiBinary, jutnto con la capacidad de multiprocesamiento, lo que permite entrenar el agente en múltiples entornos de simulación al mismo tiempo, mejorando así la eficiencia del entrenamiento y reduciendo el tiempo necesario para entrenar el agente.
Busca un equilibrio enstre la simplicidad y el rendimiento del mismo, porporcionando una implementación más sencila que otros algoritmos de optimización de políticas, como el algoritmo Trust Region Policy Optimization (TRPO), también permite una gran flexibilidad en la configuración del entorno y durante el entrenameinto del agente suele mantener una gran estabilidad, lo que lo hace adecuado para este tipo de entornos y tareas.

Las políticas disponibles del algoritmo PPO son:
\begin{itemize}
    \item \textbf{MlpPolicy:} Alias de ActorCriticPolicy.
    \item \textbf{CnnPolicy:} Alias de ActorCriticCnnPolicy.
    \item \textbf{MultiInputPolicy:} Alias de MultiInputActorCriticPolicy.
\end{itemize}
La política que se ha usado en el entrenamiento del agente es la \texttt{MultiInputPolicy} ya que es la única de las 3 que permite usar múltiples entradas, lo que es necesario para el entorno de simulación.

En este caso, el espacio de acciones es discreto, ya que el agente solo puede tomar una acción a la vez, o permitir el tráfico o denergarlo en un instante de tiempo. En cuanto al espacio de observación, que son las variables que el agente puede observar del entorno, se ha definido como un espacio de tipo Dict de \texttt{Gymnasium}, que permite definir un espacio de observación con múltiples entradas, en este caso puede observar solo la ocupación de la cola de paquetes, en un intervalo entre 0 y 1, donde 0 significa que la cola está vacía y 1 significa que la cola está llena, y el número de paquetes descartados en el último instante de tiempo, que es un número entero mayor o igual que 0.

Para calcular los descartados, el agente cuenta todos los paquetes que no han podido entrar en la cola de paquetes, ya sea porque la cola está llena o porque el agente ha decidido denegar el tráfico en ese instante de tiempo, permitiendo gestionar el tráfico de manera más eficiente y evitar la saturación de la cola de paquetes.

\subsubsection{Entrenamiento del agente}
Para la evaluación del entrenamiento del agente se han usado las métricas de recompensa media, que indica como ha ido evolucionando la recompensa media a lo largo del tiempo, y \textit{explained variance}, una métrica importante para comprobar si el agente está aprendiendo o no, que indica la calidad de las predicciones, en el que valores negativos indican que el agente obtendría mejores resultados si tomara siempre la misma acción, el mejor resultado es 1, indicando que las predicciones del agente son perfectas. ~\cite{JonathanHui2023}
En este caso estas son las estadísicas obtenidas durante el entrenamiento del agente:

\imagen{training}{Estadísticas del entrenamiento del agente}{1.0}

Como se puede observar en la imagen anterior, la recompensa media hay ido aumentando a lo largo del tiempo, mejorando con el tiempo, en cambio la explained variance ha sido más variable, aunque con un valor relativamente bueno.

\subsubsection{Evaluación del agente}
Para la evaluación del agente ya entrenado, se genera un entorno de simulación, en el que se guardan los estados de los generadores de paquetes, para posteriomente poderlo comparar con las acciones tomadas por el agente, y así poder evaluar su rendimiento.
Durante las predicciones del agente en el entorno, también se han guardado otros datos relevantes, como la ocupación de la cola en cada instante o la acción tomada.

\imagen{RewardsActions}{Estadísticas de la evaluación del agente}{1.0}

Como se puede ver en la imagen anterior, el fondo rojo indica que está en un estado de ataque, mandando paquetes de denegación de servicio hacia el agente, y el fondo blanco, lo contrario, es decir, generando tráfico benigno. En verde se puede ver el porcentaje de ocupación de la cola de paquetes, y en azul las acciones tomadas, donde 0 indica que el agente ha permitido el tráfico y 1 que lo ha denegado. Se puede apreciar que el agente ha aprendido a denegar el tráfico para que no sature la cola de paquetes, cumpliendo con el objetivo del proyecto, que es evitar la saturación de la cola de paquetes y denegar el tráfico de denegación de servicio, permitiendo así que el tráfico benigno pueda ser procesado.

\section{Desarrollo de la aplicación web}


\subsection{Despliegue de la aplicación web}
Para el despliegue de la aplicación web se ha usado un servidor virtual en Amazon Web Services (AWS) EC2, con una imagen de AMI de Amazon Linux 2023, y con un tamaño de instancia t3.medium, que dispone de los recursos necesarios para el despliegue y la ejecución de la aplicación.
Para facilitar el despliegue y la gestión de la aplicación, se ha usado Docker y Docker Compose, que permiten crear contenedores separados para cada parte de la aplicación, con cada uno sus dependencias, facilitando así su despliegue y gestión.
La arquitectura de la aplicación web se ha diseñado para que sea escalable y modular, permitiendo añadir nuevas funcionalidades en el futuro sin afectar al resto de la aplicación.
\imagen{architecture}{Arquitectura de la aplicación web}{1.0}

Los contenedores que se han creado son:
\begin{itemize}
    \item \textbf{Web:} Contiene la aplicación web desarrollada en Java con Spring Boot 3, que se encarga de servir las páginas web al usuario, y realizar toda la gestión de los modelos, junto con la visualización de información sobre el proyecto. El contenedor final generado utiliza una contrucción en multi-etapas, que en la primera etapa compila el proyecto en un archivo JAR, y en la segunda etapa se ejecuta el archivo JAR generado, permitiendo así reducir el tamaño del contenedor final, pudiendose ahorrar dependencias solo necesarias para la etapa de compilación.
    \item \textbf{API:} Contiene la API REST desarrollada en Python con Flask, que se encarga de servir los datos necesarios desde Python para la aplicación web.
    \item \textbf{Base de datos:} Contiene la base de datos MySQL, que almacena los datos necesarios para la aplicación web y la API REST. Que usa un volumen de Docker para persistir los datos, de manera que si el contenedor se elimina, los datos no se pierden.
\end{itemize}
Para la comunicación entre los contenedores se han usado dos redes de Docker, una para la comunicación entre la aplicación web y la API REST, y otra para la comunicación entre la aplicación web y la base de datos. Esto permite que cada contenedor se pueda comunicar solo y exclusivamente con los contenedores que necesita, mejorando la seguridad y el rendimiento de la aplicación.

Para el despliegue de los contenedores de Docker, se ha usado Docker Compose, que permite definir y ejecutar aplicaciones multi-contenedor a partir de la configuración de un archivo \texttt{docker-compose.yml}. Este archivo define los servicios, redes y volúmenes necesarios para la aplicación, facilitando así su despliegue y gestión.
En la configuración de Docker Compose se necesita definir las variables de entorno en el archivo \texttt{.env}.


\subsection{Seguridad del proyecto}
Uno de los pilares más importantes en el desarrollo de software es la seguridad, y en este proyecto se ha tenido en cuenta desde el principio, implementando medidas de seguridad en todos los niveles.

\subsubsection{Seguridad en el código}
Para garantizar la seguridad del código, se han seguido las siguientes prácticas:
\begin{itemize}
    \item \textbf{Control de versiones:} Se ha utilizado Git como sistema de control de versiones, lo que permite llevar un registro de todos los cambios realizados en el código, facilitando la identificación de posibles errores y vulnerabilidades.
    \item \textbf{Revisión de dependencias:} Se ha realizado una revisión automática de las dependencias utilizadas en el proyecto, para asegurarse de que no contienen vulnerabilidades conocidas. Utilizando Dependabot, que es una herramienta de GitHub que permite detectar vulnerabilidades en las dependencias del proyecto y sugerir actualizaciones para solucionarlas.

    \item \textbf{Revisión de código:} Se ha realizado una revisión de código mediante pull requests, donde se revisa el código automáticamente con GitHub Copilot, SonarQube y CodeQL analysis, para detectar posibles errores y vulnerabilidades en el código.
    También se ha utilizado herramientas de análisis estático de código de GitHub como Code Scanning y Secret Scanning, que permiten detectar posibles errores y vulnerabilidades en el código, así como secretos expuestos en el código, como contraseñas, claves API, etc, pudiendo llegar a bloquear el commit si se detecta algún secreto expuesto.

    \item \textbf{Gestión de Secretos:} Para la gestión de secretos se ha utilizado un archivo \texttt{.env} que contiene las variables de entorno necesarias para la aplicación, como las credenciales de la base de datos, las claves API, etc. Este archivo no se deberís de incluir en el repositorio de Git, para evitar que los secretos queden expuestos, debiendo ser añadido al archivo \texttt{.gitignore} para evitar que se suba al repositorio. Sin embargo, para que se pueda utilizar y probar el proyecto por parte del tribunal se ha incluido el archivo \texttt{.env}, que contienen las variables de entorno mínimas para el funcionamiento del proyecto, pero no contiene los secretos reales, como credenciales SSH para subir contenido al servidor y a GitHub y claves API de Cloudflare y de Groq. Para poder usar la funcionalidad total de la parte web se debe de acceder a la url de la aplicación web, que es \url{https://www.cesarrv.com}.
\end{itemize}

\subsubsection{Seguridad en la infraestructura}

En cuanto a la seguridad de la aplicación web, se han implementado las siguientes medidas:
\begin{itemize}
    \item \textbf{Autenticación y autorización:} Se ha implementado un sistema de autenticación con Spring Boot Security, en el que permite gestionar los usuarios y sus roles, así como proteger las rutas de la aplicación web.
    \item \textbf{Protección contra ataques CSRF:} Se ha implementado protección contra ataques CSRF implementada por Spring Boot Security.
    \item \textbf{Cifrado de datos:} Se ha guardado el hash de las contraseñas de los usuarios en la base de datos, utilizando el algoritmo BCrypt, que es un algoritmo de cifrado de contraseñas seguro y ampliamente utilizado. Esto permite proteger las contraseñas de los usuarios en caso de que la base de datos sea comprometida.
\end{itemize}

\imagen{secrets}{Secretos usados en las acciones de GitHub}{1.0}

\subsubsection{Uso de Cloudflare}
En la aplicación web se ha usado Cloudflare como servicio de protección, estadísticas, rendimiento y personalización de reglas, con el objetivo de mitigar ataques, mejorar el rendimiento, mejorar la seguridad general de la aplicación y optimización de costes.
También se ha aprovechado para comprar el dominio de \url{www.cesarrv.com} a través de Cloudflare, pudiendo así gestionar el dominio y los DNS desde la misma plataforma, facilitando la configuración de estos.
En cuanto a las solicitudes que recibe el dominio, aún redirigiendo al servidor apagado, recibe alrededor de 1000 solicitudes al día, de las cuales ninguna es legítima, ya que el dominio no está activo y el servidor donde se aloja la web no está activo. Esto indica que el dominio ha sido indexado por los motores de búsqueda y está siendo escaneado por bots maliciosos.
\imagen{cloudflare_stats}{Estadísticas de Cloudflare del dominio cesarrv.com}{1.0}
Como se puede observar en la imagen anterior, el dominio ha recibido un gran número de solicitudes. Otros datos a reseñar es que la mayoría de las solicitudes entragadas han estado guardadas en caché, lo que indica que Cloudflare ha podido servir las solicitudes sin necesidad de enviar la petición al servidor, mejorando así el rendimiento y reduciendo la carga y costos asociados del servidor.

Los principales escaneos que se han detectado son:
\begin{itemize}
    \item \textbf{Escaneos de git:} De las peticiones recibidas uno de los patrones más comunes es el de escanear el repositorio de git, buscando archivos sensibles como \texttt{.git/config}, \texttt{.git/HEAD}, \texttt{.git/index}, etc. Estos archivos pueden contener información sensible sobre la configuración del repositorio, como las ramas, los usuarios, las contraseñas, etc.
    \item \textbf{Escaneos de WordPress:} Otro patrón común detectado es el de escanear la web en busca de vulnerabilidades relacionadas con WordPress, como la búsqueda de archivos como \texttt{wp-config.php}, \texttt{wp-login.php}, \texttt{/wp-includes/}, etc. Estos archivos y directorios pueden contener información sensible sobre la configuración de WordPress, como las contraseñas, los usuarios, etc.
\end{itemize}

\subsubsection{Configuración y reglas aplicadas en Cloudflare}
Para la mitigación de los ataques descritos anteriormente, se han aplicado las siguientes reglas en Cloudflare:
\imagen{rules}{Reglas de seguridad aplicadas en Cloudflare}{1.0}
En el que la función principal de cada una es:
\begin{itemize}
    \item \textbf{Block by region:} Bloquea las solcicitudes que provienen de regiones geográficas específicas, en este caso, se han bloqueado las solicitudes que provienen de fuera de Europa, y países espacíficos dentro de Europa que se han detectado que son origen de ataques, como Rusia, Ucrania, Irlanda, etc.
    \item \textbf{Block bots:} Bloquea las solicitudes que siguen patrones comunes por bots maliciosos, centrándose en los encabezados de las peticiones y el tipo de petición que hacen.
    \item \textbf{All except Spain:} De todas las solicitudes que han pasado las anteriores reglas, si el país de la solicitud no es España, administra un desafío gestionado por Cloudflare ~\cite{CloudflareDocsTeam2025}, que requiere que el usuario resuelva un CAPTCHA para poder acceder al sitio web. Esto ayuda a filtrar las solicitudes legítimas de las maliciosas.
    \item \textbf{Too Much Requests:} Limita las solicitudes que puede hacer un mismo usuario en un periodo de tiempo determinado, en este caso, se ha configurado para que un usuario no pueda hacer más de 40 solicitudes cada 10 segundos, lo que ayuda a prevenir ataques de denegación de servicio (DoS) y ataques de escaneo por fuerza bruta.
\end{itemize}

Para la parte de seguridad frente a ataques, Cloudflare se ha configurado para evitar y/o mitigar distintos tipos comunes de ataques:
\begin{itemize}
    \item \textbf{Ataques de denegación de servicio (DoS y DDoS):} Cloudflare ayuda a mitigar ataques de denegación de servicio (DoS y DDoS) mediante la filtración de solicitudes con fines maliciosos, la limitación de solicitudes y la protección contra bots maliciosos. Además si se detecta un ataque se puede activar el modo "I'm Under Attack" ~\cite{CloudflareTeam2025}, que activa una protección adicional contra ataques DDoS, requiriendo que los usuarios resuelvan un desafío antes de acceder al sitio web.
    
    \item \textbf{Ataques XSS:} Cloudflare agrega ciertos encabezados de seguridad a las respuestas HTTP que proporcionan protección contra este tipo de ataques.

    \item \textbf{Ataques por URL:} Cloudflare hace un normalizado de las URL entrantes para que transforme los caracteres \\ en \/, juntar varias barras en una sola y aplicar la normalización de URL de la RFC 3986~\cite{CloudflareTeam2024}.    
    Todos estos pasos ayudan a prevenir ataques basados en la manipulación de URL, como los ataques de inyección de código o los ataques de redirección.
\end{itemize}
\capitulo{6}{Trabajos relacionados}

Este apartado sería parecido a un estado del arte de una tesis o tesina. En un trabajo final grado no parece obligada su presencia, aunque se puede dejar a juicio del tutor el incluir un pequeño resumen comentado de los trabajos y proyectos ya realizados en el campo del proyecto en curso. 

\capitulo{7}{Conclusiones y Líneas de trabajo futuras}

\section{Conclusiones}
El proyecto ha conseguido desarrollar un agente, que mediante el uso de aprendizaje por refuerzo es capaz de mitigar ataques de denegación de servicio en un entorno de simulación basado en dispositivos IoT.

También en la aplicación web desarrollada es importante la visualización de los datos y de la documentación del proyecto, permitiendo a los usuarios una mayor compresión del proyecto.

\section{Líneas de trabajo futuras}
Para la continuación del proyecto se han planteado varias líneas de trabajo futuras, que buscan mejorar el rendimiento del agente y la funcionalidad de la aplicación web. 

\subsection{Mejoras en el agente de aprendizaje por refuerzo}
Las principales líneas de trabajo son:
\begin{itemize}
    \item \textbf{Mejora en la recompensa:} Optimizar la función de recompensa, añadiendo más variables para que el agente tenga más información relevante y así poder aprender.
    \item \textbf{Mitigación de más ataques:} Añadir al agente la posibilidad de mitigación de más tipos de ataques como:
    \begin{itemize}
        \item Escaneo de puertos (PortScan)
        \item Escaneo de sistema operativo (OSScan)
    \end{itemize}
    \item \textbf{Desarrollar e implementar un clasificador de ataques previo al agente:} Añadiendo un clasificador de ataques que permita clasificar los ataques recibidos en diferentes categorías, para poder aplicar diferentes estrategias de mitigación según el tipo de ataque. En caso que el clasificador no pueda clasificar el ataque con suficiente confianza, se podría aplicar una estrategia de mitigación genérica por el agente.
\end{itemize}

\subsection{Mejoras en la aplicación web}
Las principales líneas de trabajo son:
\begin{itemize}
    \item \textbf{Mejora en el entorno de simulación:} Llevar la simulación a un entorno de simulación de tráfico, como GNS3, en el que se implemente el agente y se pueda ver y evaluar cual es su desempeño en un entorno más complejo.
    \item \textbf{Optimización del rendimiento:} Mejorar la eficiencia de la aplicación web, reduciendo los tiempos de carga y optimizando el uso de recursos.
    \item \textbf{Mejora de la interfaz de usuario:} Realizar cambios en el diseño y la usabilidad de la aplicación para facilitar la interacción del usuario. En cuanto a la parte técnica, se plantea la sustitución de los estilos de Bootstrap 5 por un framework de frontend moderno como \textit{React}, \textit{Vue.js} o \textit{Angular}, lo que permitiría una interfaz más dinámica, modular y fácil de mantener.
    \item \textbf{Ampliación de funcionalidades:} Incluir nuevas características en la aplicación web, como la posibilidad de ajustar los parámetros de entrenamiento del agente desde la aplicación web, permitiendo a los usuarios personalizar el comportamiento del agente según sus necesidades.
    \item \textbf{Mejora de la seguridad:} Implementar medidas adicionales para proteger la aplicación web y los datos de los usuarios, como la autenticación de dos factores y el cifrado de las comunicaciones mediante el protocolo HTTPS entre Cloudflare y el servidor de AWS.
    \item \textbf{Ampliar las pruebas:} Incrementar las pruebas de software para asegurar la calidad y estabilidad de la aplicación, haciendo énfasis en las pruebas de Selenium, para asegurar el correcto funcionamiento de la aplicación web en diferentes navegadores y dispositivos.
    \item \textbf{Mejora de la documentación del proyecto:} Ampliar y mejorar la documentación del proyecto, incluyendo guías más detalladas de usuario, manuales de instalación y uso en diferentes idiomas, mejorando las ya existentes en español y creando nuevas en inglés.
\end{itemize}



\bibliographystyle{plain}
\bibliography{bibliografia}

\end{document}
